\documentclass{article}
\usepackage{amsmath,amssymb,amsthm}

\newtheorem{lemma}{Lemma}

\begin{document}

\title{FO$^k$ Cycle-Space Invariance}
\date{}
\maketitle

\section*{Definitions}

Let $G=(V,E)$ be a graph with maximum degree $\Delta$.

Let

\[
\operatorname{tp}^k_r(v)
\]

denote the FO$^k$ $r$-type of the rooted ball $(B_r(v),v)$.

Let

\[
Z_1(B_R(v))
\]

denote the cycle space of the $R$-ball around $v$.

\section*{Lemma}

\begin{lemma}
If two vertices $u,v$ satisfy
\[
\operatorname{tp}^k_r(u)=\operatorname{tp}^k_r(v)
\]
for $r\ge R$, then

\[
Z_1(B_R(u)) \cong Z_1(B_R(v)).
\]
\end{lemma}

\begin{proof}

Equality of FO$^k$ $r$-types implies the rooted graphs
\[
(B_r(u),u) \cong (B_r(v),v)
\]
are indistinguishable in $k$-variable logic.

Since $R\le r$, the induced subgraphs

\[
B_R(u),\quad B_R(v)
\]

are isomorphic as rooted graphs.

Cycle spaces are invariants of graph isomorphism.

Therefore

\[
Z_1(B_R(u)) \cong Z_1(B_R(v)).
\]

\end{proof}

\end{document}
